\DocumentMetadata{
  lang        = en-US,
  pdfversion  = 2.0,
  pdfstandard = ua-2,
  testphase   = {tagpdf},
}

\documentclass[12pt]{book}

\tagpdfsetup{math/mathml/luamml/load=false}

\usepackage{fontspec}
\newfontfamily\cuneifont[
  Path = fonts/,
  Extension = .ttf,
  UprightFont = *-Regular
  ]{NotoSansCuneiform}
\newfontfamily\greekfont{GFS Artemisia}[
  Scale=MatchLowercase,
  Script=Greek
]
  
\usepackage{textgreek}
\usepackage{alphabeta}
\usepackage[ngerman,american]{babel}
\usepackage{csquotes}
\usepackage{caption}

\usepackage{geometry}
\geometry{
  letterpaper,
  margin=0.12\paperwidth,
  vmargin=0.10\paperheight,
  includehead,
}

\setlength{\headheight}{15pt}

\usepackage{fancyhdr}

\usepackage{graphicx}
\usepackage{array}
\usepackage{booktabs}
\usepackage{caption}
\usepackage[dvipsnames]{xcolor}
\usepackage{afterpage}
\usepackage{tocbibind}
\usepackage{subcaption}
\usepackage{float}

% \usepackage{titlesec}
% \titleformat{\chapter}[display]
%   {\normalfont\huge\bfseries}
%   {Chapter \thechapter}
%   {0.5em}
%   {\Huge}

\makeatletter
% Adjust the spacing above the chapter
\renewcommand*{\chapterheadstartvskip}{\vspace*{10pt}} 

% Redefine the 'display' look without breaking the tagging structure
\def\@makechapterhead#1{%
  {\parindent \z@ \raggedright \normalfont
    \ifnum \c@secnumdepth >\m@ne
      \if@mainmatter
        \huge\bfseries Chapter \thechapter
        \par\nobreak
        \vskip 0.5em
      \fi
    \fi
    \interlinepenalty\@M
    \Huge \bfseries #1\par\nobreak
    \vskip 40\p@
  }}
\makeatother

% Counters
\usepackage{totcount}
\usepackage{enumitem}
\usepackage{aliascnt}

\setlist[enumerate,1]{label=\textup{(\roman*)}}
\setlist[enumerate,2]{label=\textup{(\alph*)}}
\setlist[enumerate,3]{label=\textup{\arabic*.}}

\captionsetup[figure]{font=small}
\captionsetup[table]{font=small}

% Math environments
\usepackage{mathtools,amssymb,amsthm}

% Fonts and micro-typography
\usepackage{newtxtext,newtxmath}
\usepackage{microtype}
\usepackage{relsize}
\usepackage{calc}

\newtheorem{theorem}{Theorem}[section]

\newaliascnt{lemma}{theorem}
\newtheorem{lemma}[lemma]{Lemma}

\newaliascnt{proposition}{theorem}
\newtheorem{proposition}[proposition]{Proposition}

\newaliascnt{corollary}{theorem}
\newtheorem{corollary}[corollary]{Corollary}

\newtheorem{conjecture}{Conjecture}
\newtheorem{problem}[conjecture]{Problem}

\theoremstyle{definition}
\newaliascnt{definition}{theorem}
\newtheorem{definition}[definition]{Definition}

\newaliascnt{example}{theorem}
\newtheorem{example}[example]{Example}

\theoremstyle{remark}
\newaliascnt{remark}{theorem}
\newtheorem{remark}[remark]{Remark}

\newcommand{\crefnamed}[1]{\hyperref[#1]{\cref*{#1} (\nameref*{#1})}}

% TikZ and plotting
\usepackage{tikz}
\usetikzlibrary{angles,arrows.meta,babel,backgrounds,calc,decorations.pathreplacing,decorations.pathmorphing,hobby,intersections,quotes,shapes.geometric,trees}
\usepackage{pgfplots}
\pgfplotsset{compat=1.18}
\tikzset{
  every picture/.append style={line cap=round, line join=round}
}

% \title{Number Theory: A Narrative}
\title{Elementary Number Theory}
\author{Egbert Rijke}
\date{\today}

\definecolor{grayblue}{HTML}{0047AB}

\usepackage{xurl}

\usepackage{hyperref}
\hypersetup{
  unicode=true,
  colorlinks=true,
  linkcolor=grayblue,
  filecolor=grayblue,
  urlcolor=grayblue,
  citecolor=grayblue,
  pdftitle={Elementary Number Theory},
  pdfauthor  = {Egbert Rijke},
  pdfsubject = {Course Notes},
  pdfkeywords= {number theory,introduction to number theory,introduction to proof,history of number theory,Euclid,Fibonacci,Fermat,Euler,Gauss,mathematical induction,Euclidean algorithm,modular arithmetic,Pythagorean triples,infinite descent,prime numbers,Bertrand's postulate,multiplicative functions,strong divisibility sequences,Fermat's Little Theorem,Fermat's Two-Squares Theorem,polynomials,irreducible polynomials,cyclotomic polynomials,quadratic residues,quadratic reciprocity},
  hypertexnames=false,
}

\usepackage[backend=biber, style=alphabetic, sorting=nyt]{biblatex}

\DeclareLabelalphaTemplate{%
  \labelelement{%
    \field[final]{shorthand}%
    \field{label}%
    \field[strwidth=3,strside=left,ifnames=1]{labelname}%
    \field[strwidth=1,strside=left]{labelname}%
  }%
  \labelelement{%
    \field{year}%
  }%
}

\addbibresource{references.bib}

\usepackage[capitalize,english,noabbrev,nameinlink]{cleveref}

\crefname{table}{Table}{Tables}

\crefname{theorem}{Theorem}{Theorems}
\Crefname{theorem}{Theorem}{Theorems}

\aliascntresetthe{lemma}
\crefname{lemma}{Lemma}{Lemmas}
\Crefname{lemma}{Lemma}{Lemmas}

\aliascntresetthe{proposition}
\crefname{proposition}{Proposition}{Propositions}
\Crefname{proposition}{Proposition}{Propositions}

\aliascntresetthe{corollary}
\crefname{corollary}{Corollary}{Corollaries}
\Crefname{corollary}{Corollary}{Corollaries}

\aliascntresetthe{definition}
\crefname{definition}{Definition}{Definitions}
\Crefname{definition}{Definition}{Definitions}

\aliascntresetthe{example}
\crefname{example}{Example}{Examples}
\Crefname{example}{Example}{Examples}

\aliascntresetthe{remark}
\crefname{remark}{Remark}{Remarks}
\Crefname{remark}{Remark}{Remarks}

\pgfplotsset{
  lectureplot/.style={
    width=12cm, height=8cm,
    axis lines=middle,
    axis line style={-stealth},
    tick style={draw=none},
    xlabel={$x$}, ylabel={$y$},
    xmin=0, xmax=1e6,
    ymin=0, ymax=8.5e4,
    clip=false,
    legend style={
      draw=none,                               
      fill=none,
      font=\footnotesize,
      at={(0.05,0.95)},                        
      anchor=north west,
      /tikz/every even column/.append style={column sep=1em}
    },
    legend cell align=left
  }
}

\newcommand{\define}{\emph}

\newtotcounter{exercisecounter}
\newcommand{\exercise}{\stepcounter{exercisecounter}\item}

\AtEndDocument{%
  \typeout{Total number of exercises: \number\value{exercisecounter}}%
}

\newlist{exenum}{enumerate}{1}
\setlist[exenum]{label=\thechapter.\arabic*,noitemsep}

\crefname{exenumi}{Exercise}{Exercises}

\newlist{subexenum}{enumerate}{1}
\setlist[subexenum]{noitemsep,label=(\alph*),ref=\theexenumi.\alph*}
\crefname{subexenumi}{Exercise}{Exercises}

\newif\ifthischapex
\global\thischapexfalse
  
\newcommand{\exercises}{%
\par\clearpage%
\section*{Exercises}%
\addcontentsline{toc}{section}{Exercises}%
\markboth{\leftmark}{\MakeUppercase{Exercises}}%
\setcounter{exenumi}{0}%
\restartlist{exenum}
\global\thischapexfalse
}

\newcommand{\italicsubsection}[1]{%
  \subsection*{\itshape #1}%
}

\newenvironment{starter}
{%
  \italicsubsection{Starter Exercises}%
  \ifthischapex
    \begin{exenum}[resume=chapex]
  \else
    \begin{exenum}[series=chapex]
  \fi
}
{%
  \end{exenum}
  \global\thischapextrue
}

% Dutch: Routinevormende opgaven
\newenvironment{routine}
{%
  \italicsubsection{Routine-Building Exercises}%
  \ifthischapex
    \begin{exenum}[resume=chapex]
  \else
    \begin{exenum}[series=chapex]
  \fi
}
{%
\end{exenum}
\global\thischapextrue
}

\newenvironment{challenge}
{%
  \italicsubsection{Challenge Exercises}%
  \ifthischapex
    \begin{exenum}[resume=chapex]
  \else
    \begin{exenum}[series=chapex]
  \fi
}
{%
\end{exenum}
\global\thischapextrue
}

\newcommand\objectives{\italicsubsection{Learning Objectives}}

\setcounter{tocdepth}{2}

\newcommand{\Sierpinski}{Sierpiński}
\newcommand{\rpermutation}[2]{{}_{#1}P_{#2}}
\renewcommand{\pmod}[1]{\mkern5mu\left(\operatorname{mod} \mkern 6mu #1\right)}
\DeclareMathOperator{\ord}{ord}
\DeclareMathOperator{\ind}{ind}
\DeclareMathOperator*{\lcm}{lcm}
\DeclareMathOperator{\sgn}{sgn}
\DeclareMathOperator{\rad}{rad}
\newcommand{\legendre}[2]{%
  \mathchoice
    {\left(\frac{#1}{#2}\right)}% display style
    {\bigl(\tfrac{#1}{#2}\bigr)}% text style
    {\bigl(\tfrac{#1}{#2}\bigr)}% script style
    {\bigl(\tfrac{#1}{#2}\bigr)}% scriptscript style
}
\newcommand{\lres}[2]{\ensuremath{\lfloor #1 \rceil_{#2}}}
\newcommand{\abslres}[2]{\ensuremath{|\mkern-3mu \lfloor #1 \rceil\mkern-3mu |_{#2}}}
\newcommand{\complementsubset}[1]{#1^{\mathrm{c}}}
\newcommand{\Oh}{O}
\newcommand{\oh}{o}
\newcommand{\blank}{\mathord{\hspace{1pt}\text{--}\hspace{1pt}}}
\newcommand{\apart}{\mathbin{\#}}
\DeclareMathOperator{\numerator}{num}
\DeclareMathOperator{\denominator}{denom}
\newcommand{\mediant}{\mathbin{\triangledown}}
\newcommand{\bimapsto}{\mathrel{\vcenter{\offinterlineskip\ialign{##\cr$\mapsfrom$\cr\noalign{\kern0pt}$\mapsto$\cr}}}}
\newcommand{\content}{\mathrm{cont}}
\newcommand{\floorfrac}[2]{\left\lfloor\frac{#1}{#2}\right\rfloor}
\newcommand{\totient}{\phi}

\renewcommand{\headrulewidth}{0pt}
\renewcommand{\footrulewidth}{0pt}
% Header:
\fancyhead[LE]{\thepage}
\fancyhead[RE]{\itshape \leftmark}
\fancyhead[LO]{\itshape \rightmark}
\fancyhead[RO]{\thepage}
\fancyfoot{}

\fancypagestyle{noheader}{
  \fancyhf{}
  \renewcommand{\headrulewidth}{0pt}
  \setlength{\headheight}{0pt}
  \setlength{\headsep}{0pt}
}

\newif\ifshownotes
\shownotestrue     % => show notes
% \shownotesfalse      % => ignore notes

\usepackage{environ}
\NewEnviron{noteblock}{%
  \ifshownotes
    \begingroup\color{black!60}\BODY\endgroup
  \fi
}

\newcommand{\note}[1]{%
  \ifshownotes
    \begingroup\color{black!60}#1\endgroup
  \fi
}
